\chapter{The Problem in One Chapter}
\label{ch:summary}

This chapter states the $\Pc$ versus $\NP$ problem, its history, its status and its relation to randomness, in the space of a long essay. Everything in it is expanded in the chapters that follow, and the reader who already knows the field may treat it as a table of contents in prose.

\section{The question}

Some problems are easy to solve. Others are easy to check. Multiplying two numbers is easy to solve. Factoring a number is not known to be easy to solve, but a proposed factorisation is easy to check. Finding a route through every city on a map and back, shorter than a given bound, is not known to be easy to solve, but a proposed route is easy to check. Placing $N$ queens on a chessboard that already has some queens on it, so that no two attack, is not known to be easy to solve, but a proposed placement is easy to check.

Easy means solvable by a procedure whose number of steps grows no faster than a polynomial in the size of the input: $n^2$, $n^3$, $n^{10}$, but not $2^n$. The class of problems easy to solve is $\Pc$. The class of problems easy to check is $\NP$. Every problem easy to solve is easy to check, so $\Pc \subseteq \NP$. The question is whether the reverse holds: is every problem whose solutions can be checked quickly also solvable quickly? That is, is $\Pc = \NP$?

Almost everyone believes the answer is no. Nobody has proved it. It is the central open problem of computer science, one of the seven Millennium Prize Problems, and in a precise sense the question whether creativity can be mechanised: whether recognising a good answer is as hard as finding one.

\section{Where it came from}

The question has roots in the foundations of mathematics. David Hilbert\index{Hilbert, David} asked in 1928 for a mechanical procedure to decide whether any given statement of first order logic is provable, the Entscheidungsproblem. Kurt G\"odel\index{G\"odel, Kurt} showed in 1931 that no consistent, effectively axiomatized formal system containing arithmetic can prove all the true statements of arithmetic. Alonzo Church\index{Church, Alonzo} and Alan Turing\index{Turing, Alan} showed independently in 1936 that Hilbert's procedure does not exist: there is no algorithm that decides provability, because there is no algorithm that decides whether a given algorithm halts. Turing's proof introduced the abstract machine that bears his name, and the machine is the model of computation in which $\Pc$ and $\NP$ are defined.

Turing showed what cannot be computed at all. The question of what can be computed efficiently came later. G\"odel asked it in a letter to von Neumann\index{von Neumann, John} in 1956, observing that if a machine could find proofs of length $n$ in time proportional to $n$ or $n^2$, mathematics could be mechanised in practice despite his own theorem. Hartmanis and Stearns defined time complexity classes and proved the time hierarchy theorem in 1965. Cobham\index{Cobham, Alan} and Edmonds\index{Edmonds, Jack} proposed polynomial time as the criterion of feasibility in the same year. Cook\index{Cook Levin theorem} in 1971 and Levin in 1973 defined $\NP$ completeness and proved that satisfiability of Boolean formulas is complete; Karp\index{Karp, Richard} in 1972 showed that twenty one natural problems are complete; and the list has grown since to thousands, including the queens completion problem in 2017.

\section{What is known}

Remarkably little, and most of it is about what will not work. Three barriers are known. Baker, Gill and Solovay\index{Baker, Gill and Solovay} showed in 1975 that any proof that relativizes, that goes through unchanged when all machines are given access to the same oracle, cannot resolve the question, because there are oracles relative to which $\Pc = \NP$ and oracles relative to which $\Pc \neq \NP$. Razborov\index{Razborov, Alexander} and Rudich showed in 1994 that natural proofs, those that establish a lower bound by exhibiting a property of hard functions that is both easy to recognise and shared by many functions, cannot work unless strong pseudorandom functions do not exist, which would contradict standard cryptographic assumptions. Aaronson and Wigderson\index{Wigderson, Avi} showed in 2008 that algebrizing proofs, which extend relativization to allow the oracle to be queried on a low degree extension, also cannot work. The techniques that proved the known separations, diagonalisation for the time hierarchy and combinatorial arguments for restricted circuit classes, fall on one side or another of these barriers.

What is known on the positive side is a web of implications. If $\Pc = \NP$ then the entire polynomial hierarchy collapses to $\Pc$, cryptography based on one way functions is impossible, and every problem in $\NP$ is solved in polynomial time by a single universal algorithm, Levin's. If $\Pc \neq \NP$ then by Ladner\index{Ladner's theorem}'s theorem there are problems in $\NP$ neither in $\Pc$ nor $\NP$ complete. If $\NP \subseteq \Ppoly$ then the hierarchy collapses to its second level, by Karp and Lipton. And the lower bounds that have been proved, for monotone circuits by Razborov, for constant depth circuits by Furst, Saxe, Sipser\index{Sipser, Michael} and by H{\aa}stad, and for the class $\mathrm{ACC}^0$ against $\mathrm{NEXP}$ by Williams\index{Williams, Ryan} in 2011, are each a long way from what $\Pc \neq \NP$ would need, which is a superpolynomial lower bound on the size of general circuits for some problem in $\NP$. The best known lower bounds for explicit functions against general circuits are linear, between $3n$ and $5n$ depending on the gate basis.

\section{Randomness}

A separate question, more tractable and equally fundamental, is whether a random coin helps. A randomized algorithm may flip coins and may err with small probability. The class $\RP$ is the problems solvable in polynomial time by an algorithm that never says yes wrongly and says yes to a true instance with probability at least a half; $\BPP$ allows errors on both sides. Primality was in $\coRP$ from 1977 and was found to be in $\Pc$ in 2002 by Agrawal\index{Agrawal, Kayal and Saxena}, Kayal and Saxena, the model of a successful derandomization. Polynomial identity testing is in $\coRP$ and not known to be in $\Pc$.

The belief that $\Pc = \BPP$, that randomness is a convenience and not a resource, rests on the hardness versus randomness theorems of Nisan and Wigderson and of Impagliazzo\index{Impagliazzo and Wigderson} and Wigderson: if some problem solvable in exponential time requires exponentially large circuits, then there are pseudorandom generators that fool circuits of any fixed polynomial size, with the seed length depending on that size, which suffices to derandomize $\BPP$, and every randomized algorithm can be run on the outputs of the generator instead of on true random bits. The hypothesis is a circuit lower bound of the kind nobody can prove, and Kabanets and Impagliazzo showed that this is no accident: derandomizing even polynomial identity testing would prove a circuit lower bound. So $\Pc = \RP$ is believed, is implied by believed hypotheses, and is blocked by the same wall as $\Pc \neq \NP$.

\section{Counting}

Harder than deciding is counting. The number of solutions of a problem in $\NP$ is a function in the class $\sharpP$ of Valiant\index{Valiant, Leslie}, and counting is $\sharpP$ complete for the permanent, for satisfying assignments, for perfect matchings, and for queens completions. Toda\index{Toda's theorem} showed that a single counting oracle decides every problem in the polynomial hierarchy. But approximate counting is easier: Stockmeyer\index{Stockmeyer's theorem} showed it can be done with an $\NP$ oracle and random bits, and Jerrum, Valiant and Vazirani\index{Jerrum, Valiant and Vazirani} showed that for self reducible\index{self reducibility} problems approximate counting is polynomially equivalent to approximately uniform sampling. The algorithms that count in practice, from the permanent to the volume of a convex body, are Markov chains that sample, and their analysis is the analysis of mixing times.

\section{This book}

The queens problem is our laboratory. It is a counting problem, $\sharpP$ complete in its completion form, with an energy, the number of attacking pairs, that makes it a Boltzmann\index{Boltzmann distribution} system in the sense of statistical physics. We show that counting its solutions, and more generally computing any expectation over an exponentially large space, can be written by a Fubini argument as the integral over the unit interval of a monotone step function, the quantile function of the integrand. This is the Lorenz\index{Lorenz curve} identity, and the Monte Carlo methods built on it, vertical likelihood, nested sampling and split sampling, walk down the staircase from the top. The walk is polynomial in every respect except one: the production of a sample from a level set of the energy, which is a Markov chain mixing problem.

The thesis is that this is also the right picture for $\Pc$ versus $\RP$. A randomized algorithm has a quantile function\index{quantile function} with a single step, and deciding membership in an $\RP$ language is evaluating that function at the midpoint. In these coordinates the hitting set formulation of $\Pc = \RP$ (a sufficient condition for it, not known to be equivalent) asks that a dense level set, of density at least a half, be hit by a deterministic polynomial size set of points; the thesis is that $\NP \neq \RP$ corresponds to a sparse level set, of density $e^{-\Omega(N)}$ and shattered into exponentially many components, that cannot be sampled by any polynomial chain, a correspondence we do not prove. The two statements are about the same staircase at different heights, and the book is an attempt to make that observation precise, to test it on the queens, and to say what remains to be proved.

\section{What the book finds}

Since the thesis was first stated the queens have answered some of the questions put to them, and a reader should know the answers before the argument begins. On unconditioned boards the solutions are not isolated: for $N \le 13$ all but a handful lie in one connected component of the level set with at most one attack, bridged by near solutions, and the swap chain on that component has autocorrelation time growing like $N^{1.8}$ to $N = 32$. The component is held together by paths through near solutions, some of which lie eight swaps from the nearest solution, not by direct repairs, and a proof of polynomial mixing on it reduces, by a theorem of Chapter~\ref{ch:mixingproof}, to a single conjecture, the inverse polynomial expansion of a family of finite graphs. That conjecture, if true, gives polynomial mixing at the bottom rung of the ladder, which is one ingredient of a polynomial time Monte Carlo algorithm for $Q(N)$; such an algorithm would also need that the component carries nearly all solutions and polynomial mixing on the upper rungs, neither of which we prove.

On conditioned boards the picture splits. Boards cut out of a solution keep the giant component to $N = 17$. Boards conditioned by random queens do not: at $N = 12$ about half of those with two to four fixed queens have their completions scattered over several components of the bottom level set, and an adversary can scatter them completely. But on every such board examined the completions reconnect two attacks up and the level set itself reconnects two to four attacks up, and a theorem of Chapter~\ref{ch:open} says the walk may skip the broken rungs at a price equal to the ratio of the ground state to the first connected level, about $10^{-3}$ at $N = 12$, some thirty times smaller than the ratio across an unbroken rung. The Monte Carlo half of the program is thereby sharpened, given mixing above it and inverse polynomial rung ratios, to one quantity: the shattering depth of a board, the number of attacks one must allow before its level set is connected. On random boards it is two to four at $N = 12$, three on most boards, and three or four at $N = 14$ and $16$, with four more common at the larger size, so the data do not yet show whether it grows with $N$ there. On the boards that encode hard formulas we conjecture it grows with $N$, and we have not measured it.

Two methods that might have been thought to avoid the question do not. Rubinstein's cross entropy method with a product proposal cannot beat crude Monte Carlo on the queens by more than a factor of about $2$ for $7 \le N \le 12$, by a three line theorem of Chapter~\ref{ch:ce}, and his cloning method is the vertical walk under another name. A proposal tilted toward Simkin's queenon does not accelerate the chain. And the derandomization half of the book, which gives it its title, remains a restatement: $\Pc = \RP$ follows if a level set of density one half can be hit without a coin, and the quantile coordinate makes the claim clean without making it easier.
